% 数学结论及完整证明已移入正文相应公式之后。


% 数学结论及完整证明已移入正文相应公式之后。


% 数学结论及完整证明已移入正文相应公式之后。


% 数学结论及完整证明已移入正文相应公式之后。


% 数学结论及完整证明已移入正文相应公式之后。


\begin{derivation}[从\eqnrefs{eq:green_discrete_mode_spectral_representation}到\eqnrefs{eq:lossless_mode_sum_to_green_fdt_zeroT}]
这一谱分解要求介质线性、无损，并满足使 Maxwell 算符自伴的边界条件； 吸收介质不具有同样的离散实频正交谱，应使用连续储库描述。 自伴性来自分部积分：

\[
\langle\mathbf u,\hat\Theta\mathbf v\rangle
=\int\dd^3 r\,
[\nabla\times\boldsymbol\varepsilon_r^{-1/2}\mathbf u]^*\cdot
\boldsymbol\mu_r^{-1}
[\nabla\times\boldsymbol\varepsilon_r^{-1/2}\mathbf v],
\]
在所选边界条件下交换 $\mathbf u,\mathbf v$ 不产生边界项。$\{\mathbf u_\lambda\}$ 的完备性通过
$\mathbf u_\lambda=\boldsymbol\varepsilon_r^{1/2}\mathbf f_\lambda$
转化为
\[
\sum_\lambda
\boldsymbol\varepsilon_r(\mathbf r)\mathbf f_\lambda(\mathbf r)
\otimes\mathbf f_\lambda^*(\mathbf r')
=\mathbf I\delta^{(3)}(\mathbf r-\mathbf r').
\]
对每个 $\mathbf f_\lambda$，Maxwell 预解式分母为
\[
\left[\nabla\times\boldsymbol\mu_r^{-1}\nabla\times
-\frac{(\omega+\ii0^+)^2}{c^2}\boldsymbol\varepsilon_r\right]\mathbf f_\lambda
=\frac{\omega_\lambda^2-(\omega+\ii0^+)^2}{c^2}
\boldsymbol\varepsilon_r\mathbf f_\lambda.
\]
将完备关系逐模除以这一本征值即得相应的 Green 谱展开。再用
\[
\operatorname{Im}\frac1{\omega_\lambda^2-(\omega+\ii0^+)^2}
=\frac{\pi}{2\omega}\delta(\omega-\omega_\lambda)
\quad(\omega>0)
\]
得到模求和恒等式。最后把单量子电场矩阵元
\[
\mathbf E_\lambda^{(+)}\otimes\mathbf E_\lambda^{(-)}
=\frac{\hbar\omega_\lambda}{2\varepsilon_0}
\mathbf f_\lambda\otimes\mathbf f_\lambda^*
\]
乘以在壳 $\delta$ 并求和，使用 $1/(\varepsilon_0c^2)=\mu_0$，得到零温 FDT 形式。
\end{derivation}

% 数学结论及完整证明已移入正文相应公式之后。


\begin{derivation}[从\eqnrefs{eq:global-continuum-normalization-delta}到\eqnrefs{eq:general_mode_coupling_operator}]
先用单位归一周期盒离散化连续动量谱，再由单个动量格点的体积恢复连续核 $\mathcal G^{(\lambda)}_{\mathbf p's',\mathbf ps}$。盒体积 $V$ 只承担正则化作用；在 $V\to\infty$ 时，它必须在“态密度 $\times$ 矩阵元”的组合中消去。取 $V=L_xL_yL_z$，允许动量为

\begin{equation*}
\mathbf p_{\mathbf n}
=2\pi\hbar\left(\frac{n_x}{L_x},\frac{n_y}{L_y},\frac{n_z}{L_z}\right),
\end{equation*}
盒态归一化为
\begin{equation*}
{}_V\!\langle\mathbf p_{\mathbf n'},s'|\mathbf p_{\mathbf n},s\rangle_V
=\delta_{\mathbf n'\mathbf n}\delta_{s's}.
\end{equation*}
当 $V\to\infty$ 时，每个动量格点占据体积
\begin{equation*}
\Delta^3p=\frac{(2\pi\hbar)^3}{V},
\end{equation*}
所以
\begin{equation*}
\sum_{\mathbf n}F(\mathbf p_{\mathbf n})
\longrightarrow
\frac{V}{(2\pi\hbar)^3}\int\dd^3 p\,F(\mathbf p).
\end{equation*}
与连续约定
$\langle\mathbf p',s'|\mathbf p,s\rangle=\delta_{s's}\delta^{(3)}(\mathbf p'-\mathbf p)$
相容的态变换是
\begin{equation*}
|\mathbf p_{\mathbf n},s\rangle_V
=\sqrt{\Delta^3p}\,|\mathbf p_{\mathbf n},s\rangle.
\end{equation*}
因此盒中一量子顶点矩阵元
\begin{equation*}
\mathcal G^{(\lambda,V)}_{\mathbf n's',\mathbf ns}
=\frac1\hbar{}_V\!\langle\mathbf p_{\mathbf n'},s'|
\delta h_\lambda|\mathbf p_{\mathbf n},s\rangle_V
\end{equation*}
与连续核之间满足
\begin{equation*}
\mathcal G^{(\lambda,V)}_{\mathbf n's',\mathbf ns}
=\Delta^3p\,
\mathcal G^{(\lambda)}_{\mathbf p_{\mathbf n'}s',\mathbf p_{\mathbf n}s}
\end{equation*}
（若初末态使用不同离散化体积，相应平方根因子分别保留；最终物理组合不变）。离散 Schr\"odinger 方程中的耦合和为
\begin{equation*}
\sum_{\mathbf n}
\mathcal G^{(\lambda,V)}_{\mathbf n'\mathbf n}C_{\mathbf n}.
\end{equation*}
代入上式并用 $\sum_{\mathbf n}\Delta^3p\to\int\dd^3 p$，得到
\begin{equation*}
\int\dd^3 p\,
\mathcal G^{(\lambda)}_{\mathbf p'\mathbf p}C(\mathbf p).
\end{equation*}
所以盒体积在“态密度 $\times$ 盒矩阵元”这一组合中完全消去；连续核本身不是一个普通频率，而只有
\begin{equation*}
[\mathcal G_{\mathbf p'\mathbf p}\,\dd^3 p]={\rm s^{-1}}
\end{equation*}
才具有跃迁率量纲。这也解释了为什么连续谱中的 $\delta$ 归一化、态密度与耦合矩阵元不能拆开单独赋予物理概率。
\end{derivation}

% 数学结论及完整证明已移入正文相应公式之后。


% 数学结论及完整证明已移入正文相应公式之后。


% 数学结论及完整证明已移入正文相应公式之后。


% 数学结论及完整证明已移入正文相应公式之后。


% 数学结论及完整证明已移入正文相应公式之后。


% 数学结论及完整证明已移入正文相应公式之后。


% 数学结论及完整证明已移入正文相应公式之后。


\begin{derivation}[从\eqnrefs{eq:universal_free_electron_boson_hamiltonian}到\eqnrefs{eq:universal_boson_finite_time_emission}]
对矩形作用窗完成时间积分并取模平方，就得到正文\eqnrefs{eq:universal_boson_finite_time_emission} 的 $\sinc^2$ 谱核。 在相互作用绘景中，电子算符与模式算符分别按自由 Hamiltonian 演化。 对电子本征态，

\begin{align*}
{}_V\!\langle f_e|\delta\hat h_\lambda^\dagger(t)|i_e\rangle_V
&=\ee^{\ii(E_f-E_i)t/\hbar}
{}_V\!\langle f_e|\delta\hat h_\lambda^\dagger|i_e\rangle_V,\\
\langle n_\lambda+1|\hat a_\lambda^\dagger(t)|n_\lambda\rangle
&=\ee^{+\ii\omega_\lambda t}\sqrt{n_\lambda+1}.
\end{align*}
因此发射末态 $|f_e;\mathbf n+\mathbf e_\lambda\rangle_V$ 与初态之间的瞬时矩阵元为
\begin{align*}
\langle f_e;\mathbf n+\mathbf e_\lambda|\hat H_I(t)|i_e;\mathbf n\rangle
&=M_{fi,\lambda}^{\rm em}\sqrt{n_\lambda+1}
\exp\!\left\{\ii\left[\frac{E_f-E_i}{\hbar}+\omega_\lambda\right]t\right\}\\
&=M_{fi,\lambda}^{\rm em}\sqrt{n_\lambda+1}
\ee^{\ii\Delta_{fi}^{\rm em}t}.
\end{align*}
Dyson 级数的一阶项于是给
\[
C_f^{(1)}(T)
=-\frac{\ii}{\hbar}M_{fi,\lambda}^{\rm em}\sqrt{n_\lambda+1}
\int_0^T\dd t\,\ee^{\ii\Delta_{fi}^{\rm em}t}.
\]
时间积分可以写成
\[
\int_0^T\dd t\,\ee^{\ii\Delta t}
=T\ee^{\ii\Delta T/2}\sinc(\Delta T/2),
\]
故
\[
C_f^{(1)}(T)
=-\ii\mathcal G_{fi,\lambda}^{\rm em}\sqrt{n_\lambda+1}
T\ee^{\ii\Delta_{fi}^{\rm em}T/2}
\sinc(\Delta_{fi}^{\rm em}T/2).
\]
取模平方时中间相位 $\exp(\ii\Delta T/2)$ 消失，得到\eqnrefs{eq:universal_boson_finite_time_emission}。其中 $n_\lambda+1$ 直接来自 $|\langle n_\lambda+1|a_\lambda^\dagger|n_\lambda\rangle|^2$；即使 $n_\lambda=0$，产生算符仍有非零矩阵元。吸收支路的末态为 $|f_e;\mathbf n-\mathbf e_\lambda\rangle_V$。此时
\begin{align*}
{}_V\!\langle f_e|\delta\hat h_\lambda(t)|i_e\rangle_V
&=\ee^{\ii(E_f-E_i)t/\hbar}M_{fi,\lambda}^{\rm abs},\\
\langle n_\lambda-1|\hat a_\lambda(t)|n_\lambda\rangle
&=\ee^{-\ii\omega_\lambda t}\sqrt{n_\lambda},
\end{align*}
所以总相位的角频率失谐为
\[
 \Delta_{fi}^{\rm abs}
 =\frac{E_f-E_i}{\hbar}-\omega_\lambda
 =\frac{E_f-E_i-\hbar\omega_\lambda}{\hbar}.
\]
因此
\begin{align*}
 C_{f,\rm abs}^{(1)}(T)
 &=-\ii\mathcal G_{fi,\lambda}^{\rm abs}\sqrt{n_\lambda}
 T\,\ee^{\ii\Delta_{fi}^{\rm abs}T/2}
 \sinc(\Delta_{fi}^{\rm abs}T/2),\\
 P_{i\to f}^{\rm abs}(T)
 &=n_\lambda|\mathcal G_{fi,\lambda}^{\rm abs}|^2T^2
 \sinc^2(\Delta_{fi}^{\rm abs}T/2).
\end{align*}
这与发射支路的 $n_\lambda+1$ 和 $\Delta_{fi}^{\rm em}$ 构成两条独立的有限时间结果。
\end{derivation}

% 数学结论及完整证明已移入正文相应公式之后。


\begin{derivation}[从\eqnrefs{eq:general_free_electron_golden_rule}到\eqnrefs{eq:generalized_cherenkov_condition}]
保留 Hessian 项可得到首个反冲修正； 只有再删去这一项，才得到正文\eqnrefs{eq:generalized_cherenkov_condition}。 取
$\delta\mathbf p=-\hbar(\boldsymbol\kappa+\mathbf G)$。 多变量 Taylor 展开给

\[
E(\mathbf p_i+\delta\mathbf p)
=E(\mathbf p_i)+\mathbf v_i\cdot\delta\mathbf p
+\frac12\delta\mathbf p^{\mathsf T}H_E(\mathbf p_i)\delta\mathbf p
+O(\|\delta\mathbf p\|^3).
\]
代入精确能量守恒
$E(\mathbf p_i)-E(\mathbf p_i+\delta\mathbf p)=\hbar\omega_{\nu\boldsymbol\kappa}$。一次项满足
\[
-\mathbf v_i\cdot\delta\mathbf p
=\hbar\,\mathbf v_i\cdot(\boldsymbol\kappa+\mathbf G),
\]
二次项为
\[
-\frac12\delta\mathbf p^{\mathsf T}H_E\delta\mathbf p
=-\frac{\hbar^2}{2}(\boldsymbol\kappa+\mathbf G)^{\mathsf T}
H_E(\mathbf p_i)(\boldsymbol\kappa+\mathbf G).
\]
两边除以 $\hbar$，得到
\[
\omega_{\nu\boldsymbol\kappa}
=\mathbf v_i\cdot(\boldsymbol\kappa+\mathbf G)
-\frac{\hbar}{2}(\boldsymbol\kappa+\mathbf G)^{\mathsf T}
H_E(\mathbf p_i)(\boldsymbol\kappa+\mathbf G)
+O(\hbar^2\|\boldsymbol\kappa+\mathbf G\|^3).
\]
对自由相对论电子，
\[
H_{E,ij}(\mathbf p_i)
=\frac{c^2}{E_i}\delta_{ij}
-\frac{c^4p_{i,i}p_{i,j}}{E_i^3}.
\]
若电子主要沿 $z$ 传播，定义入射 Lorentz 因子 $\gamma_{{\rm L},i}\equiv E_i/(m_{\mathrm e} c^2)$，并记 $\mathbf K\equiv\boldsymbol\kappa+\mathbf G$，则
\[
\mathbf K^{\mathsf T}H_E\mathbf K
=\frac{K_\perp^2}{m_\perp}+\frac{K_z^2}{m_\parallel},
\qquad
m_\perp=\gamma_{{\rm L},i}m_{\mathrm e},
\quad
m_\parallel=\gamma_{{\rm L},i}^3m_{\mathrm e}.
\]
因此首个反冲修正具有直接的纵、横曲率分解：
\[
\omega_{\nu\boldsymbol\kappa}
=v_iK_z-
\frac{\hbar}{2}
\left(\frac{K_\perp^2}{m_\perp}+\frac{K_z^2}{m_\parallel}\right)+\cdots.
\]
当曲率项在实际模式带宽内小于实验分辨率或有限时间线宽时，最低阶同步关系~\eqnrefs{eq:generalized_cherenkov_condition} 才成为可用近似。
\end{derivation}

\begin{derivation}[从\eqnrefs{eq:general_sequential_two_boson_dyson_amplitude}到\eqnrefs{eq:two_boson_resolvent_amplitude_general}]
取时间无关耦合在有限散射窗内绝热开关，并写发射顶角的相互作用绘景矩阵元
\[
\langle m,n+1|\hat V_+(t_1)|i,n\rangle
=\sqrt{n+1}\,V^{\rm em}_{mi}
\exp\!\left[
\frac{\ii(E_m-E_i+\hbar\omega_0)t_1}{\hbar}
\right],
\]
\[
\langle f,n+2|\hat V_+(t_2)|m,n+1\rangle
=\sqrt{n+2}\,V^{\rm em}_{fm}
\exp\!\left[
\frac{\ii(E_f-E_m+\hbar\omega_0)t_2}{\hbar}
\right].
\]
代入二阶 Dyson 振幅，定义
\[
\Omega_1=\frac{E_m-E_i+\hbar\omega_0}{\hbar},
\qquad
\Omega_2=\frac{E_f-E_m+\hbar\omega_0}{\hbar},
\]
则时间核为
\[
I_T(\Omega_2,\Omega_1)
=\int\dd t_2\int^{t_2}\dd t_1\,
\ee^{\ii\Omega_2t_2}\ee^{\ii\Omega_1t_1}.
\]
时间排序区域 $t_1<t_2$ 最适合改用“总时间”和“两次顶角的时间间隔”作坐标：
\begin{equation*}
 t\equiv t_2,
 \qquad
 \tau\equiv t_2-t_1\ge0,
 \qquad
 t_1=t-\tau,
 \qquad
 \dd t_2\,\dd t_1=\dd t\,\dd\tau.
\end{equation*}
指数因子于是精确分解成
\begin{align*}
 \ee^{\ii\Omega_2t_2}\ee^{\ii\Omega_1t_1}
 &=\ee^{\ii\Omega_2t}
   \ee^{\ii\Omega_1(t-\tau)}\\
 &=\ee^{\ii(\Omega_1+\Omega_2)t}
   \ee^{-\ii\Omega_1\tau}.
\end{align*}
$t$ 积分只检查整个过程的总能量守恒，而 $\tau\ge0$ 记录第一顶角之后中间态能够传播的持续时间；两个积分因此分别编码整体守恒律与中间态传播。

为把半轴积分定义成分布，在相对时间上加入收敛因子 $\ee^{-\epsilon\tau}$，$\epsilon>0$，最后取 $\epsilon\to0^+$：
\begin{align*}
 I_\epsilon(\Omega_2,\Omega_1)
 &=\int_{-\infty}^{\infty}\dd t\,
   \ee^{\ii(\Omega_1+\Omega_2)t}
   \int_0^\infty\dd\tau\,
   \ee^{-(\epsilon+\ii\Omega_1)\tau}\\
 &=2\pi\delta(\Omega_1+\Omega_2)
   \frac1{\epsilon+\ii\Omega_1}\\
 &=2\pi\delta(\Omega_1+\Omega_2)
   \frac{-\ii}{\Omega_1-\ii\epsilon}.
\end{align*}
总时间部分给
\begin{equation*}
2\pi\delta(\Omega_1+\Omega_2)
=2\pi\hbar\,
\delta(E_f-E_i+2\hbar\omega_0).
\end{equation*}
另一方面，二阶 Dyson 前因子是 $(-\ii/\hbar)^2=-1/\hbar^2$。因此相对时间积分与 Dyson 前因子的乘积为
\begin{align*}
 -\frac1{\hbar^2}
 \frac{-\ii}{\Omega_1-\ii0^+}
 &=\frac{\ii}{\hbar^2}
 \frac{\hbar}{E_m-E_i+\hbar\omega_0-\ii0^+}\\
 &=-\frac{\ii}{\hbar}
 \frac1{E_i-E_m-\hbar\omega_0+\ii0^+}.
\end{align*}
于是完整二阶 $S$ 矩阵元具有标准的形式
\begin{align*}
 \langle f,n+2|S^{(2)}|i,n\rangle
 ={}&-2\pi\ii\,
 \delta(E_f-E_i+2\hbar\omega_0)\\
 &\times\sqrt{(n+1)(n+2)}
 \sum_m
 \frac{V^{\rm em}_{fm}V^{\rm em}_{mi}}
 {E_i-E_m-\hbar\omega_0+\ii0^+}.
\end{align*}
去掉散射矩阵中共同的 $-2\pi\ii\delta(E_f-E_i+2\hbar\omega_0)$ 后，约化二阶振幅正是
\[
\mathcal M_{fi}^{(2,\,2{\rm em})}
=\sqrt{(n+1)(n+2)}
\sum_m
\frac{V^{\rm em}_{fm}V^{\rm em}_{mi}}
{E_i-E_m-\hbar\omega_0+\ii0^+}.
\]
若电子中间谱连续，$\sum_m$ 展开为对完整动量、自旋和其它量子数的积分。利用
\[
\frac1{x+\ii0^+}
=\mathcal P\frac1x-\ii\pi\delta(x),
\]
可见中间分母的主值部分对应离壳虚路径，而 $\delta(x)$ 部分表示中间态本身可以在壳传播。后一情形若在有限脉冲带宽内占显著权重，就应恢复显式级联通道，而不应继续把全过程压成局域二阶有效顶角。
\end{derivation}

% 数学结论及完整证明已移入正文相应公式之后。


\begin{derivation}[从\eqnrefs{eq:quantum_transition_current_green_rate}到\eqnrefs{eq:dipole_green_decay_rate}]
连续性方程并不能唯一决定 $\mathbf j_{fi}(\mathbf r)$ 的每一点，因此不能把它简单“等同”为一个 $\delta$ 型电流。真正需要的是它的空间积分。取时间依赖 $\ee^{-\ii\omega_{if}t}$，则
\[
-\ii\omega_{if}\rho_{fi}
+\boldsymbol\nabla\cdot\mathbf j_{fi}=0.
\]
将第 $k$ 个分量乘以 $r_k$ 并对空间积分；对局域束缚跃迁，边界通量在无穷远消失，所以分部积分给
\begin{align*}
-\ii\omega_{if}
\int\dd^3r\,r_k\rho_{fi}(\mathbf r)
&=-\int\dd^3r\,r_k\partial_i j_{fi,i}(\mathbf r)\\
&=\int\dd^3r\,j_{fi,k}(\mathbf r).
\end{align*}
因此
\[
\boxed{
\int\dd^3r\,\mathbf j_{fi}(\mathbf r)
=-\ii\omega_{if}\mathbf d_{fi}.}
\]
这条积分恒等式才是偶极极限所需的输入。

令跃迁电流主要支撑在以 $\mathbf r_0$ 为中心、半径 $a$ 的区域内。把 Green 张量写成
\[
\operatorname{Im}\mathbf G^R(\mathbf r,\mathbf r';\omega)
=\operatorname{Im}\mathbf G^R(\mathbf r_0,\mathbf r_0;\omega)
+\Delta\mathbf G(\mathbf r,\mathbf r';\omega).
\]
代入一般跃迁率中的双重空间积分，常数项因上述积分恒等式因子化为
\begin{align*}
&\int\dd^3r\dd^3r'\,
\mathbf j_{fi}^*(\mathbf r)\cdot
\operatorname{Im}\mathbf G^R(\mathbf r_0,\mathbf r_0;\omega_{if})
\cdot\mathbf j_{fi}(\mathbf r')\\
&\qquad=\omega_{if}^2
\mathbf d_{fi}^*\cdot
\operatorname{Im}\mathbf G^R(\mathbf r_0,\mathbf r_0;\omega_{if})
\cdot\mathbf d_{fi}.
\end{align*}
其余全部多极修正收进
\[
R_{\rm mp}
\equiv
\int\dd^3r\dd^3r'\,
\mathbf j_{fi}^*(\mathbf r)\cdot
\Delta\mathbf G(\mathbf r,\mathbf r';\omega_{if})
\cdot\mathbf j_{fi}(\mathbf r').
\]
若工作区域内 Green 张量一阶可微，则均值定理给
\[
\|\Delta\mathbf G\|
\le
2a\,
\sup_{\mathcal B_a(\mathbf r_0)^2}
\|\boldsymbol\nabla_{(\mathbf r,\mathbf r')}\operatorname{Im}\mathbf G^R\|,
\]
从而
\[
|R_{\rm mp}|
\le
2a\,
\sup\|\boldsymbol\nabla\operatorname{Im}\mathbf G^R\|
\left(\int\dd^3r\,\|\mathbf j_{fi}(\mathbf r)\|\right)^2.
\]
当这个余项相对于偶极主项可忽略时，代回正文\eqnrefs{eq:quantum_transition_current_green_rate} 就得到\eqnrefs{eq:dipole_green_decay_rate}。进一步取真空
\[
\operatorname{Im}\mathbf G_0^R(\mathbf r_0,\mathbf r_0;\omega)
=\frac{\omega}{6\pi c}\mathbb I
\]
便恢复
\[
\Gamma_0
=\frac{\omega^3|\mathbf d_{fi}|^2}{3\pi\varepsilon_0\hbar c^3},
\]
从而同时检查了 Green 张量归一化与 SI 系数。
\end{derivation}

% 数学结论及完整证明已移入正文相应公式之后。


% 数学结论及完整证明已移入正文相应公式之后。


\begin{derivation}[从\eqnrefs{eq:nonlocal_permittivity_kernel}到\eqnrefs{eq:macroscopic_qed_fdt_green_tensor}]
先把空间指标、坐标积分全部压缩成算符记号。固定实频 $\omega>0$，Maxwell 算符与其伴随为
\[
\mathsf L
=\nabla\times\nabla\times
-\frac{\omega^2}{c^2}\boldsymbol\varepsilon,
\qquad
\mathsf L^\dagger
=\nabla\times\nabla\times
-\frac{\omega^2}{c^2}\boldsymbol\varepsilon^\dagger,
\]
并定义
\[
\mathsf G^R=\mathsf L^{-1},
\qquad
\mathsf G^{R\dagger}=(\mathsf L^\dagger)^{-1}.
\]
这里 $\boldsymbol\varepsilon$ 可以是空间非局域积分算符；因此以下代数不使用
$\varepsilon(\mathbf r)\delta(\mathbf r-\mathbf r')$ 的局域假设。

对任意可逆算符 $A,B$，
\[
A^{-1}-B^{-1}=A^{-1}(B-A)B^{-1}.
\]
取 $A=\mathsf L$、$B=\mathsf L^\dagger$，得到
\begin{align*}
\mathsf G^R-\mathsf G^{R\dagger}
&=\mathsf G^R(\mathsf L^\dagger-\mathsf L)\mathsf G^{R\dagger}\\
&=\frac{\omega^2}{c^2}
\mathsf G^R(\boldsymbol\varepsilon-\boldsymbol\varepsilon^\dagger)
\mathsf G^{R\dagger}\\
&=2\ii\frac{\omega^2}{c^2}
\mathsf G^R\boldsymbol\varepsilon_I\mathsf G^{R\dagger}.
\end{align*}
两边除以 $2\ii$：
\[
\operatorname{Im}\mathsf G^R
=\frac{\omega^2}{c^2}
\mathsf G^R\boldsymbol\varepsilon_I\mathsf G^{R\dagger}.
\]
展开成坐标核即为正文\eqnrefs{eq:nonlocal_green_optical_identity}，其中两个中间坐标积分来自两个积分算符的乘法。

下一步量子化耗散核。因为被动性保证
$\boldsymbol\varepsilon_I\succeq0$，谱定理允许取
\[
\boldsymbol\varepsilon_I
=\mathsf R_\varepsilon\mathsf R_\varepsilon^\dagger.
\]
正文噪声电流可写成算符形式
\[
|\hat j_N(\omega)\rangle
=\omega\sqrt{\frac{\hbar\varepsilon_0}{\pi}}
\mathsf R_\varepsilon(\omega)|\hat f(\omega)\rangle.
\]
储库真空满足
\[
\langle0|
|\hat f(\omega)\rangle\langle\hat f(\omega')|
|0\rangle
=\mathsf I\,\delta(\omega-\omega'),
\]
故
\begin{align*}
\langle0|
|\hat j_N(\omega)\rangle\langle\hat j_N(\omega')|
|0\rangle
&=\frac{\hbar\varepsilon_0}{\pi}\omega^2
\mathsf R_\varepsilon\mathsf R_\varepsilon^\dagger
\delta(\omega-\omega')\\
&=\frac{\hbar\varepsilon_0}{\pi}\omega^2
\boldsymbol\varepsilon_I
\delta(\omega-\omega').
\end{align*}

再用
\[
|\hat E(\omega)\rangle
=\ii\mu_0\omega\mathsf G^R(\omega)|\hat j_N(\omega)\rangle
\]
得到
\begin{align*}
\langle0|
|\hat E(\omega)\rangle\langle\hat E(\omega')|
|0\rangle
&=\frac{\hbar\mu_0^2\varepsilon_0}{\pi}\omega^4
\mathsf G^R\boldsymbol\varepsilon_I\mathsf G^{R\dagger}
\delta(\omega-\omega')\\
&=\frac{\hbar\mu_0}{\pi}\omega^2
\operatorname{Im}\mathsf G^R
\delta(\omega-\omega'),
\end{align*}
其中第二步使用
$\mu_0\varepsilon_0=c^{-2}$ 与前面证明的 Green 光学恒等式。展开空间和 Cartesian 指标便得到
\[
\boxed{
\langle0|
\hat E_i(\mathbf r,\omega)
\hat E_j^\dagger(\mathbf r',\omega')
|0\rangle
=\frac{\hbar\mu_0}{\pi}\omega^2
[\operatorname{Im}\mathbf G^R(\mathbf r,\mathbf r';\omega)]_{ij}
\delta(\omega-\omega').}
\]

若介质空间局域，取
\[
\boldsymbol\varepsilon_I(\mathbf s,\mathbf s';\omega)
=\boldsymbol\varepsilon_I(\mathbf s,\omega)
\delta^{(3)}(\mathbf s-\mathbf s'),
\]
双重积分立即缩成常见单积分公式。若介质互易，还可进一步使用
$\mathbf G^R(\mathbf r,\mathbf r')=\mathbf G^{R\mathsf T}(\mathbf r',\mathbf r)$；但这两条都不是上面推导成立所需的前提。对于完全无损但开放的结构，$\boldsymbol\varepsilon_I=0$ 只说明体耗散消失；无穷远辐射边界产生的表面项仍必须保留，或通过把散射通道显式纳入储库来恢复同一谱恒等式。
\end{derivation}

% 数学结论及完整证明已移入正文相应公式之后。


\begin{derivation}[从\eqnrefs{eq:macroscopic_qed_fdt_green_tensor}到\eqnrefs{eq:quantum_transition_current_green_rate}]
先把正文使用的跃迁连续性方程本身推出来。设电子在同一个静态外场 Hamiltonian
\begin{equation*}
\hat h_D
=c\boldsymbol\alpha\cdot[-\ii\hbar\boldsymbol\nabla-q_{\mathrm e}\mathbf A(\mathbf r)]
+\beta m_{\mathrm e} c^2+q_{\mathrm e}\phi(\mathbf r)
\end{equation*}
中有两个定态
\begin{equation*}
\Psi_n(\mathbf r,t)=\psi_n(\mathbf r)\ee^{-\ii E_nt/\hbar},
\qquad n=i,f.
\end{equation*}
定义交叉电荷密度与电流
\begin{equation*}
\rho_{fi}(\mathbf r,t)
=q_{\mathrm e}\Psi_f^\dagger\Psi_i,
\qquad
\mathbf j_{fi}(\mathbf r,t)
=q_{\mathrm e} c\Psi_f^\dagger\boldsymbol\alpha\Psi_i.
\end{equation*}
由两组 Dirac 方程
$\ii\hbar\partial_t\Psi_n=\hat h_D\Psi_n$，
\begin{align*}
\partial_t\rho_{fi}
&=\frac{\ii q_{\mathrm e}}{\hbar}
\bigl[(\hat h_D\Psi_f)^\dagger\Psi_i
-\Psi_f^\dagger\hat h_D\Psi_i\bigr].
\end{align*}
质量项与标势项在方括号中逐项相消；矢势项也因左右两边含同一个乘法算符而相消。只剩空间导数：
\begin{align*}
\partial_t\rho_{fi}
&=-q_{\mathrm e} c\left[
(\boldsymbol\nabla\Psi_f^\dagger)\cdot\boldsymbol\alpha\Psi_i
+\Psi_f^\dagger\boldsymbol\alpha\cdot\boldsymbol\nabla\Psi_i
\right]\\
&=-\boldsymbol\nabla\cdot\mathbf j_{fi}.
\end{align*}
所以交叉双线性与普通概率流一样满足精确连续性方程
\begin{equation*}
\partial_t\rho_{fi}+\boldsymbol\nabla\cdot\mathbf j_{fi}=0.
\end{equation*}
又因为
\begin{equation*}
\rho_{fi}(\mathbf r,t)
=\rho_{fi}(\mathbf r)\ee^{-\ii(E_i-E_f)t/\hbar}
=\rho_{fi}(\mathbf r)\ee^{-\ii\omega_{if}t},
\end{equation*}
故
\begin{equation*}
\boxed{
\boldsymbol\nabla\cdot\mathbf j_{fi}(\mathbf r)
=\ii\omega_{if}\rho_{fi}(\mathbf r)}.
\end{equation*}
这说明正文中的跃迁流守恒不是对经典轨迹的假设，而是同一 Dirac Hamiltonian 两个本征态之间的局域恒等式。

现在创建一个频率 $\omega_\lambda>0$ 的环境量子。记相应负频率场矩阵元为
$\phi_\lambda^{(-)},\mathbf A_\lambda^{(-)},\mathbf E_\lambda^{(-)}$；它们带时间因子
$\ee^{+\ii\omega_\lambda t}$，所以
\begin{equation*}
\mathbf E_\lambda^{(-)}
=-\boldsymbol\nabla\phi_\lambda^{(-)}
-\ii\omega_\lambda\mathbf A_\lambda^{(-)}.
\end{equation*}
最小耦合的一阶矩阵元为
\begin{equation*}
M_{fi,\lambda}
=\int\dd^3 r\,
[\rho_{fi}\phi_\lambda^{(-)}
-\mathbf j_{fi}\cdot\mathbf A_\lambda^{(-)}].
\end{equation*}
对第一项所需的梯度做分部积分；若跃迁流在无穷远足够快衰减，边界项为零：
\begin{align*}
\int\dd^3 r\,\mathbf j_{fi}\cdot\mathbf E_\lambda^{(-)}
&=-\int\dd^3r\,\mathbf j_{fi}\cdot\boldsymbol\nabla\phi_\lambda^{(-)}
-\ii\omega_\lambda\int\dd^3r\,\mathbf j_{fi}\cdot\mathbf A_\lambda^{(-)}\\
&=\int\dd^3r\,
(\boldsymbol\nabla\cdot\mathbf j_{fi})\phi_\lambda^{(-)}
-\ii\omega_\lambda\int\dd^3r\,\mathbf j_{fi}\cdot\mathbf A_\lambda^{(-)}\\
&=\ii\omega_{if}\int\dd^3r\,\rho_{fi}\phi_\lambda^{(-)}
-\ii\omega_\lambda\int\dd^3r\,\mathbf j_{fi}\cdot\mathbf A_\lambda^{(-)}.
\end{align*}
Fermi 黄金规则中的能量守恒分布为
\begin{equation*}
\delta(E_f+\hbar\omega_\lambda-E_i)
=\frac1\hbar\delta(\omega_\lambda-\omega_{if}),
\end{equation*}
所以真正贡献速率的支撑上有 $\omega_\lambda=\omega_{if}$。因此
\begin{equation*}
\boxed{
M_{fi,\lambda}
=\frac{1}{\ii\omega_{if}}
\int\dd^3 r\,
\mathbf j_{fi}(\mathbf r)\cdot
\mathbf E_\lambda^{(-)}(\mathbf r)}.
\end{equation*}
标势与矢势至此已经由守恒跃迁流组合成物理电场；这也是最终结果不依赖势的规范选择的直接原因。

零温下，环境初态是真空，只允许创建一个环境量子。黄金规则给
\begin{align*}
\Gamma_{i\to f}^{(0)}
&=\frac{2\pi}{\hbar}
\sum_\lambda |M_{fi,\lambda}|^2
\delta(E_f+\hbar\omega_\lambda-E_i)\\
&=\frac{2\pi}{\hbar^2\omega_{if}^2}
\int\dd^3 r\,\dd^3 r'\,
\mathbf j_{fi}^*(\mathbf r)\cdot
\left[\sum_\lambda
\mathbf E_\lambda^{(+)}(\mathbf r)\otimes
\mathbf E_\lambda^{(-)}(\mathbf r')
\delta(\omega_{if}-\omega_\lambda)\right]
\cdot\mathbf j_{fi}(\mathbf r').
\end{align*}
宏观 QED 的涨落--耗散谱恒等式\eqnrefs{eq:macroscopic_qed_fdt_green_tensor} 恰好把方括号中的“所有未分辨环境模式之和”写成
\begin{equation*}
\sum_\lambda
\mathbf E_\lambda^{(+)}(\mathbf r)\otimes
\mathbf E_\lambda^{(-)}(\mathbf r')
\delta(\omega-\omega_\lambda)
=\frac{\hbar\mu_0}{\pi}\omega^2
\operatorname{Im}\mathbf G^R(\mathbf r,\mathbf r';\omega).
\end{equation*}
在 $\omega=\omega_{if}$ 处代入，$\omega_{if}^2$ 与前面的
$1/\omega_{if}^2$ 精确抵消：
\begin{align*}
\Gamma_{i\to f}^{(0)}
&=\frac{2\pi}{\hbar^2\omega_{if}^2}
\frac{\hbar\mu_0}{\pi}\omega_{if}^2
\iint\dd^3r\,\dd^3r'\,
\mathbf j_{fi}^*(\mathbf r)\cdot
\operatorname{Im}\mathbf G^R(\mathbf r,\mathbf r';\omega_{if})
\cdot\mathbf j_{fi}(\mathbf r')\\
&=\boxed{
\frac{2\mu_0}{\hbar}
\iint\dd^3r\,\dd^3r'\,
\mathbf j_{fi}^*(\mathbf r)\cdot
\operatorname{Im}\mathbf G^R(\mathbf r,\mathbf r';\omega_{if})
\cdot\mathbf j_{fi}(\mathbf r')}.
\end{align*}
这就是正文\eqnrefs{eq:quantum_transition_current_green_rate}。电子态的全部信息留在 $\mathbf j_{fi}$，而未分辨环境模的全部信息被 Green 张量谱函数统一求和。
\end{derivation}

% 数学结论及完整证明已移入正文相应公式之后。


\begin{derivation}[从\eqnrefs{eq:prescribed_current_green_loss_master}到\eqnrefs{eq:green_tensor_eels_master}]
EELS 中真正需要证明的是：为什么耗散只由推迟 Green 张量的反厄米部分，也就是 \(\operatorname{Im}\mathbf G^R\)，决定。对给定经典源，电子转移给环境的能量定义为诱导场对源做功的负值，
\begin{equation*}
W_{\rm loss}
=-\int\dd t\,\dd^3r\,
\mathbf J(\mathbf r,t)\cdot\mathbf E_{\rm ind}(\mathbf r,t).
\end{equation*}
采用正文的 Fourier 约定，实电流满足
\(\widetilde{\mathbf J}(\mathbf r,-\omega)
=\widetilde{\mathbf J}^{*}(\mathbf r,\omega)\)。把正负频率配对，得到
\begin{equation*}
W_{\rm loss}
=-\int_0^\infty\frac{\dd\omega}{\pi}
\operatorname{Re}
\int\dd^3r\,
\widetilde{\mathbf J}^{*}(\mathbf r,\omega)
\cdot\widetilde{\mathbf E}_{\rm ind}(\mathbf r,\omega).
\end{equation*}
线性 Maxwell 响应给出
\begin{equation*}
\widetilde{\mathbf E}_{\rm ind}(\mathbf r,\omega)
=\ii\mu_0\omega
\int\dd^3r'\,
\mathbf G^R(\mathbf r,\mathbf r';\omega)
\cdot\widetilde{\mathbf J}(\mathbf r',\omega).
\end{equation*}
定义二次型
\begin{equation*}
X(\omega)
\equiv
\iint\dd^3r\,\dd^3r'\,
\widetilde{\mathbf J}^{*}(\mathbf r,\omega)
\cdot\mathbf G^R(\mathbf r,\mathbf r';\omega)
\cdot\widetilde{\mathbf J}(\mathbf r',\omega).
\end{equation*}
则
\begin{equation*}
-\operatorname{Re}[\ii X]
=\operatorname{Im}X.
\end{equation*}
而对任意核算符，
\begin{equation*}
\operatorname{Im}\mathbf G^R
\equiv\frac{\mathbf G^R-(\mathbf G^R)^\dagger}{2\ii},
\end{equation*}
所以二次型的虚部严格等于
\begin{equation*}
\operatorname{Im}X
=\iint\dd^3r\,\dd^3r'\,
\widetilde{\mathbf J}^{*}(\mathbf r,\omega)
\cdot\operatorname{Im}\mathbf G^R(\mathbf r,\mathbf r';\omega)
\cdot\widetilde{\mathbf J}(\mathbf r',\omega).
\end{equation*}
因此
\begin{equation*}
W_{\rm loss}
=\int_0^\infty\dd\omega\,
\frac{\mu_0\omega}{\pi}
\iint\dd^3r\,\dd^3r'\,
\widetilde{\mathbf J}^{*}
\cdot\operatorname{Im}\mathbf G^R
\cdot\widetilde{\mathbf J}.
\end{equation*}
因此 \(\operatorname{Re}\mathbf G^R\) 只改变反应场的保守相位，而净能量损失只取决于反厄米谱权重。把每个 \(\hbar\omega\) 的能量损失除以量子能量，便得到正文的损失数谱一般表达式。

对点电子轨迹
\begin{equation*}
\mathbf J_e(\mathbf r,t)
=q_{\mathrm e}\mathbf v(t)\,
\delta^{(3)}[\mathbf r-\mathbf r_e(t)],
\end{equation*}
其频域电流为
\begin{equation*}
\widetilde{\mathbf J}_e(\mathbf r,\omega)
=q_{\mathrm e}\int\dd t\,\mathbf v(t)
\delta^{(3)}[\mathbf r-\mathbf r_e(t)]
\ee^{\ii\omega t}.
\end{equation*}
把左右两个电流分别代入二次型，\(\mathbf r\) 与 \(\mathbf r'\) 积分由两个空间 \(\delta\) 函数精确完成：
\begin{equation*}
\frac{\dd P}{\dd\omega}
=\frac{\mu_0q_{\mathrm e}^2}{\pi\hbar}
\int\dd t\,\dd t'\,
\ee^{\ii\omega(t'-t)}
\mathbf v(t)\cdot
\operatorname{Im}\mathbf G^R[\mathbf r_e(t),\mathbf r_e(t');\omega]
\cdot\mathbf v(t').
\end{equation*}
利用 \(\operatorname{Im}\mathbf G^R\) 的厄米核性质交换 \(t,t'\) 可得到正文采用的相位顺序；再令 \(\mathbf v(t)=\mathbf v_0\) 就得到正文\eqnrefs{eq:green_tensor_eels_master}。因此 EELS 轨迹公式不是另一个经验模型，而是一般线性响应二次型在点电子源上的精确约化。
\end{derivation}

\begin{derivation}[从\eqnrefs{eq:general_free_electron_mode_coupling_definition}到\eqnrefs{eq:mode-emission-absorption-finite-time-dp68}]
对单一模式 $\lambda$，相互作用绘景中与一量子交换有关的部分写成
\[
\hat H_I(t)
=\hbar\mathcal G_{fi}^{(\lambda)}
|f\rangle\langle i|\hat a_\lambda^\dagger
 \ee^{\ii\Delta_{\rm em}t}
+\hbar\mathcal G_{fi}^{(\lambda)}
|f\rangle\langle i|\hat a_\lambda
 \ee^{\ii\Delta_{\rm abs}t}
+\mathrm{H.c.},
\]
其中
\[
\Delta_{\rm em}
=\frac{E_f-E_i+\hbar\omega_\lambda}{\hbar},
\qquad
\Delta_{\rm abs}
=\frac{E_f-E_i-\hbar\omega_\lambda}{\hbar}.
\]
一阶 Dyson 振幅对发射过程给
\begin{align*}
\mathcal A_{\rm em}^{(1)}
&=-\frac{\ii}{\hbar}
\int_{-T/2}^{T/2}\dd t\,
\langle f,n_\lambda+1|\hat H_I(t)|i,n_\lambda\rangle\\
&=-\ii\mathcal G_{fi}^{(\lambda)}\sqrt{n_\lambda+1}
\,T\,\sinc\!\left(\frac{\Delta_{\rm em}T}{2}\right).
\end{align*}
这里 $\sqrt{n_\lambda+1}$ 完全来自
$\hat a_\lambda^\dagger|n_\lambda\rangle
=\sqrt{n_\lambda+1}|n_\lambda+1\rangle$。因此
\[
P_{\rm em}^{(1)}(T)
=|\mathcal G_{fi}^{(\lambda)}|^2(n_\lambda+1)T^2
\sinc^2\!\left(\frac{\Delta_{\rm em}T}{2}\right).
\]
同理，吸收过程使用
$\hat a_\lambda|n_\lambda\rangle=\sqrt{n_\lambda}|n_\lambda-1\rangle$，得到
\[
P_{\rm abs}^{(1)}(T)
=|\mathcal G_{fi}^{(\lambda)}|^2n_\lambda T^2
\sinc^2\!\left(\frac{\Delta_{\rm abs}T}{2}\right).
\]
这就是正文\eqnrefs{eq:mode-emission-absorption-finite-time-dp68}。

有限时间核本身已经给出能量分辨率。令
\[
K_T(\Delta)
\equiv\frac{T}{2\pi}\sinc^2\!\left(\frac{\Delta T}{2}\right).
\]
变量代换 $u=\Delta T/2$ 给
\[
\int_{-\infty}^{\infty}\dd\Delta\,K_T(\Delta)
=\frac1\pi\int_{-\infty}^{\infty}\dd u\,\sinc^2u=1.
\]
因此 $K_T$ 是单位面积、宽度随 $T^{-1}$ 缩小的近似恒等核；只有在末态连续且其余因子在该宽度内变化缓慢时，$K_T$ 才在分布极限收缩为 $\delta(\Delta)$。所以 $n_\lambda+1$ 的真空项来自玻色代数，而能量守恒 $\delta$ 来自长时间极限，两者是完全不同的数学来源。
\end{derivation}

% 数学结论及完整证明已移入正文相应公式之后。


% 数学结论及完整证明已移入正文相应公式之后。


% 数学结论及完整证明已移入正文相应公式之后。


% 数学结论及完整证明已移入正文相应公式之后。


% 数学结论及完整证明已移入正文相应公式之后。


% 数学结论及完整证明已移入正文相应公式之后。


% 数学结论及完整证明已移入正文相应公式之后。


% 数学结论及完整证明已移入正文相应公式之后。


% 数学结论及完整证明已移入正文相应公式之后。


% 数学结论及完整证明已移入正文相应公式之后。


% 数学结论及完整证明已移入正文相应公式之后。


% 数学结论及完整证明已移入正文相应公式之后。


% 数学结论及完整证明已移入正文相应公式之后。


\begin{derivation}[直接积分连续谱，而不预先代入能量守恒]
写 $T^2\sinc^2(\Delta T/2)=\int_0^T\dd t\int_0^T\dd t'\,\ee^{\ii\Delta(t-t')}$，并使用 Lorentz 线形的 Fourier 积分
\begin{equation*}
 \int_{-\infty}^{\infty}\dd\Delta\,
 \frac{J_0\ee^{\ii\Delta u}}{1+\Delta^2\tau_c^2}
 =\frac{\pi J_0}{\tau_c}\ee^{-|u|/\tau_c}.
\end{equation*}
两时间正方形区域中，相同时间差 $u>0$ 对应的线段长度为 $T-u$，故
\begin{align*}
 P^{(2)}(T)&=\frac{2\pi J_0}{\tau_c}
 \int_0^T(T-u)\ee^{-u/\tau_c}\dd u\\
 &=2\pi J_0[T-\tau_c(1-\ee^{-T/\tau_c})].
\end{align*}
因此恒定速率是有限时间积分的渐近结果。谱阈值、窄峰或离散末态改变了被积分的 $\mathcal J_{\rm tr}$ 时，应先重算这个积分，再判断是否存在黄金律区。
\end{derivation}
